% 08_conclusion.tex —— 优缺点与结论
\section{Strengths and Weaknesses}

The main strength is consistency. Problems~1--4 share one symbol set, one
forward model and one economic convention. The pressure redistribution of
Problem~2 is validated against the three-phase simulation of Problem~1; the
reduced production model of Problem~3 is calibrated to the training groups of
Problem~4; and the economic evaluation follows Appendix~B exactly, including
the reference state, the daily discounting and the fracturing-only baseline.
The robustness metrics are separated into feasibility, conditional
profitability and joint frequency, which keeps an infeasible candidate from
being reported as merely uneconomic.

The limitations are equally explicit. The model is one-dimensional normal to
the fracture, so pressure loss along the fracture and lateral interference are
neglected. Natural fractures in Case~3 are represented by an equivalent
permeability rather than a dual-porosity model. Water production is not
supplied, so the produced-water term uses a constant water--oil ratio. The
damage variable is inactive in the nominal pressure range, and the effect of
injected volume and CO$_2$ fraction on production is not identifiable from the
supplied data; the zero-exponent convention and the positive-exponent control
are both reported.

\section{Conclusion}

For Problem~1, the single-phase benchmark shows approximately second-order
spatial and first-order temporal convergence. Table~\ref{tbl:verification}
gives mass-balance errors for the tested flow settings. For Problem~2, the
Case~3 numerical 5\% pressure radius vanishes between 3 and 4~d at
300~m$^3$/d and between 6 and 8~d at 800~m$^3$/d. Both analytical
comparisons give earlier loss of the threshold in these tests.

For Problem~3, Case~1 has no feasible point in the prescreen. Three low-rate
three-phase runs exceed the pressure limit before the planned volume is injected. Case~2 is
feasible but uneconomic: the compromise candidates give NPV between $-164$
and $-87\times10^4$~CNY, with a joint feasibility--profitability frequency of
zero. Case~3 is the only robustly viable scenario: NPV ranges from $+365$ to
$+671\times10^4$~CNY, all 200 samples per candidate pass the reduced prescreen and have positive NPV. The compromise soak times are $8$--$24$~d under the calibrated
production model.

The engineering implications are conditional. Further pressure tests would be needed to establish a viable injection window
for Case~1. Case~2 remains uneconomic for the three selected candidates at
the baseline oil price. Case~3 supports the
lowest-cost operating point under the zero-exponent convention; if a positive
volume or CO$_2$-fraction benefit is later demonstrated, the preferred injected
volume moves from $1000$ to about $1668$~m$^3$. The absolute NPV values should
be read together with the uncalibrated water-production and baseline terms.
